\chapter{Results and Discussion}
\section{Checking DE on test functions}
All five methods were given the same number of evaluations. On the ten-dimensional Sphere function, DE's median error is \SphereDE, which is essentially zero. This confirms that our from-scratch DE works correctly.

\ArtifactTable{tab_benchmarks}{Final error on each test function: mean and standard deviation over all runs. Lower is better.}{tab:benchmarks}

\ArtifactFigurePair{fig_benchmark_convergence}{Benchmark convergence across all runs}{Best error found so far against the number of evaluations (median run with middle-50\% band).}{fig:benchconvergence}{fig_benchmark_box}{Distribution of final benchmark errors}{Spread of final errors over repeated runs for each method.}{fig:benchbox}

As expected, no single method wins on every function. DE converges reliably on Sphere and has the best median on the curved Rosenbrock valley. The self-adapting jDE does best on the trap-filled Rastrigin and Griewank functions. Errors close to $10^{-15}$ simply mean the method has reached the optimum.

\ArtifactFigureShort[.9]{fig_de_rastrigin_2d}{DE population snapshots on Rastrigin}{Snapshots of the DE population on the two-dimensional Rastrigin function over successive generations.}{fig:rastrigin}

\section{How the DE search progressed}
The DE search ran \NumEvals\ trials. The population's average fitness rose quickly toward the best, and later generations made smaller refinements. Diversity shrank as the population gathered around good settings, but never reached zero, and children were still replacing parents in the last generation. The dashed baseline line in \Cref{fig:deconv} uses the same quick training setup, so the two are directly comparable.

\ArtifactFigurePair{fig_de_convergence}{DE proxy-fitness progression}{Best and average validation fitness per generation; the shaded band shows the spread of the population.}{fig:deconv}{fig_de_diversity}{Population diversity and replacement success}{Population diversity and the share of children that replaced their parents in each generation.}{fig:diversity}

\ArtifactFigurePair{fig_search_landscape}{Candidate fitness projected onto two genes}{Every trial plotted by learning rate and fine-tuning depth, coloured by fitness.}{fig:landscape}{fig_parallel_coords}{Parallel coordinates of candidate recipes}{All trial recipes, one line each, across the seven genes.}{fig:parallel}

\ArtifactFigureShort{fig_gene_importance}{Gene--fitness associations}{How strongly each gene was linked with higher fitness across all trials.}{fig:importance}

\section{Chosen recipes and training curves}
The full-size retraining step picked the final DE recipe from its best few candidates. This recipe retrains the whole network, keeps the baseline's head size, and adds label smoothing and weight decay, with no extra RandAugment. Random search chose a recipe that retrains part of the network, uses a wider head, and applies stronger augmentation. \Cref{fig:importance} agrees: trials that retrained more of the network and used less dropout tended to score higher.

Both searched recipes reach a clearly higher validation macro-F1 than the baseline. Their loss values look higher only because label smoothing changes how the loss is measured; accuracy and macro-F1 are the fair comparison.

\ArtifactTable{tab_best_config}{Baseline and the final recipes chosen by DE and random search.}{tab:recipes}

\ArtifactFigureShort{fig_learning_curves}{Final-training learning curves}{Training curves of the final models (mean and spread over seeds). Compare models by macro-F1 or accuracy, since label smoothing changes the loss scale.}{fig:learning}

\section{Final test results}
The DE recipe raises mean test accuracy from \BaseAcc\% to \DEAcc\% and macro-F1 from \BaseFone\% to \DEFone\%. Its results are also very stable across seeds, varying by only \DEStd\ percentage points. In other words, DE removes \ErrorReduction\% of the baseline's mistakes. Random search reaches the same \RSAcc\% accuracy and \RSFone\% macro-F1, which confirms that the improvement comes from tuning the training recipe.

\ArtifactTable{tab_final_metrics}{Final test results: mean and standard deviation over three training seeds.}{tab:final}

\ArtifactFigureShort{fig_before_after_bars}{Final held-out metric comparison}{Test metrics of the baseline, DE, and random-search models, with variation across seeds. Higher is better.}{fig:hero}

\ArtifactFigurePair{fig_confusion_side_by_side}{Paired seed-zero confusion matrices}{Confusion matrices for the seed-zero models. Each row is the true class; off-diagonal cells are mistakes.}{fig:confusion}{fig_per_class_f1}{Per-class F1 comparison}{Average F1 for each class, sorted by how much DE improved over the baseline.}{fig:perclass}

\ArtifactFigureShort[.8]{fig_roc}{Classwise and macro ROC curves}{ROC curves for each class and their average, from the test predictions.}{fig:roc}

\textbf{Every class improves under DE.} Early blight and Target Spot gain the most, although Early blight remains the hardest class. The baseline's most common mistakes were Target Spot confused with spider mites, Early blight with Late blight, and Target Spot with healthy leaves. DE removes the Target Spot mistakes. Its few remaining errors are mostly between the two blight diseases, whose leaf spots can look alike.

\section{Is the improvement real?}
For the seed-zero models, the 95\% confidence range for accuracy is \BaseAccCI\ for the baseline and \DEAccCI\ for DE; for macro-F1 it is \BaseFoneCI\ and \DEFoneCI. The ranges do not overlap. The McNemar test gives $p=\McNemarP$, which confirms that DE's improvement over the baseline is not due to chance.

\ArtifactTable{tab_confidence_intervals}{95\% bootstrap confidence intervals for the seed-zero models.}{tab:intervals}

\section{Robustness, confidence, and Grad-CAM}
On the corruption test (\Cref{tab:robust,fig:robust}), the DE model handles blur and JPEG compression better than the baseline, but is more affected by noise and strong lighting changes. Its average corruption accuracy is \DERobustAcc\% against \BaseRobustAcc\% for the baseline. Its calibration error is \DEECE\% against \BaseECE\%; the reliability diagram shows that DE is mostly \emph{under}-confident, meaning it is right more often than its confidence suggests, which fits with the softer targets from label smoothing. Since the search optimised only clean accuracy, this shows that clean accuracy alone does not guarantee robustness. Adding robustness to the search goal is a natural next step (Chapter~11).

\ArtifactTable{tab_robustness}{Calibration error, average corruption accuracy, and speed of the seed-zero models.}{tab:robust}

\ArtifactFigurePair{fig_robustness}{Corruption stress test}{Accuracy under each type and strength of image damage.}{fig:robust}{fig_reliability}{Confidence reliability diagrams}{Reliability diagrams: how model confidence compares with actual correctness.}{fig:reliability}

\ArtifactFigureShort[1]{fig_gradcam_grid}{Six held-out Grad-CAM cases}{Grad-CAM heat maps for the first test image of six classes, for the seed-zero baseline and DE models. Brighter areas influenced the prediction more.}{fig:gradcam}

Both models classify all six example leaves correctly. The baseline's heat maps sometimes focus on leaf edges or the background, while DE's maps usually cover more of the leaf itself. Chapter~9 looks at three of these leaves in detail. The repository also contains a prediction gallery and a feature-space plot as extra material.
